% ==============================================================================
% secoes/06_conclusao.tex - SEÇÃO VI: CONCLUSÃO
% Tema: Federico Faggin e o Primeiro Microprocessador Monolítico (Intel 4004)
% ==============================================================================

\section{Conclusão}
\label{sec:conclusao}

O Intel 4004 superou a tirania dos números (\cref{eq:tirania_numeros}), segundo a qual a confiabilidade cai conforme cresce a quantidade de componentes interligados, e atendeu à previsão de Moore ao reunir funções aritméticas, lógicas e de controle em uma única pastilha de silício \cite{morton1958, moore1965}. Essa consolidação eliminou centenas de conexões externas entre módulos discretos que elevavam o custo de montagem e a taxa de falhas dos computadores. A análise reunida nas \cref{sec:fundamentacao,sec:resultados} comprova que arquiteturas lógicas só se tornam fabricáveis quando integradas à física de dispositivos e ao \textit{layout}. Ao substituir a porta metálica de alumínio pelo eletrodo de silício policristalino na tecnologia de porta autoalinhada (SGT), o processo contornou os limites capacitivos e de tensão de limiar dos transistores PMOS anteriores \cite{faggin1970}. Sobre essa base tecnológica, as soluções implementadas por Federico Faggin e sintetizadas na \cref{tab:inovacoes_faggin}, centradas nos contatos enterrados, no \textit{bootstrapping} dinâmico e no roteamento ortogonal em duas camadas, levaram a CPU à produção comercial \cite{faggin1996, faggin2021}.

No plano industrial e econômico, o conjunto MCS-4 consolidou a transição dos circuitos dedicados sob encomenda, precursores dos modernos ASICs (\textit{Application-Specific Integrated Circuits}, chips projetados para uma única função fixa), para componentes programáveis de catálogo apoiados em memórias de leitura (ROM e EPROM) e de acesso aleatório (RAM) (\cref{subsec:reorganizacao_economica}) \cite{faggin1972fourchip}. Ao transferir a lógica dos produtos de ligações físicas rígidas para instruções gravadas em memória, esse modelo reduziu os custos iniciais de projeto e consolidou os sistemas embarcados no controle de equipamentos médicos, comerciais e industriais. A partir dessa herança técnica, os processadores de 8 bits Intel 8008, Intel 8080 e Zilog Z80 impulsionaram a automação fabril, ilustrada pelo robô elétrico ASEA IRB 6, e abriram caminho para a primeira geração da microcomputação pessoal (\cref{subsec:impactos_sociedade}) \cite{noyce1981, chmz80panel2007, bls1982futurework}.

Na trajetória física posterior dos semicondutores, as restrições de velocidade e dissipação dos transistores PMOS orientaram a passagem para o processo NMOS, baseado em elétrons mais velozes, e depois para a lógica complementar CMOS, que reduziu o consumo em repouso e sustentou o escalamento proporcional de Dennard por mais de três décadas (\cref{subsec:limitacoes_pmos_nmos,fig:transicao_mos_dennard}) \cite{dennard1974, mead1980}. Quando a miniaturização geométrica atingiu a barreira de potência (limite de calor dissipado por área) e esbarrou em efeitos quânticos (vazamento de elétrons através de isolantes ultrafinos por tunelamento), o encolhimento isolado dos transistores deixou de gerar ganhos proporcionais de desempenho \cite{cavin2012}. Diante dessas barreiras físicas, a adoção da co-otimização entre projeto e tecnologia (DTCO, \textit{Design-Technology Co-Optimization}, prática em que físicos de materiais, engenheiros de fábrica e arquitetos projetam em conjunto desde o início) reafirma a atualidade do método integrado aplicado por Faggin no Intel 4004.
